\begin{tabular}{p{2.9cm}p{3.0cm}p{3.0cm}p{5.5cm}}
\hline
Grandezza & V5 & V6 & nota \\
\hline
Rapporto $\omega_2/\omega_1$ & 0{,}119667646 & $45/376=0{,}119680851$ & il valore V5 non e' $8{,}3565$ (cio\`e $0{,}11966732$) ma $1/8{,}35648$: non ricostruibile dai dati dichiarati \\
$I_2/I_1$ & 8{,}353816 & $376/45=8{,}355556$ & scarto $-2{,}1\times10^{-4}$; la chiusura richiede la massa dell'App.~A \\
$r_{\rm eff}$ & 239{,}93 nm & 60{,}26 / 259{,}34 nm & volume di materiale / raggio idrodinamico: sono due grandezze diverse \\
Massa & $5{,}73\times10^{-18}$ kg & $8{,}21\times10^{-18}$ kg & la V5 contava i soli fili; la V6 aggiunge il nucleo (29\,\% del totale) \\
Forza $F$ & $3{,}08\times10^{-17}$ N & $4{,}52\times10^{-17}$ N & $+47$\,\%: raggi di filo 3,2 / 7,5 / 4,5 nm, non 1,25 nm \\
Coppia $\tau$ & $2{,}61\times10^{-24}$ N\,m & $3{,}31\times10^{-24}$ N\,m & $+27\%$: $P_{\rm ass}$ e non $P_{\rm int}$, con le sezioni corrette \\
$Pe_r$ & $6{,}10\times10^{-4}$ & $7{,}72\times10^{-4}$ & identit\`a $\tau/k_BT$; la V5 riportava anche $8{,}8\times10^{-3}$ con un $\beta$ diverso \\
\hline
\end{tabular}
